\chapter{Musculoskeletal System}

\label{ch:ch54_musculoskeletal}

\section{What does musculoskeletal imaging cover?}
The musculoskeletal (MSK) system includes bones, joints, muscles, tendons, ligaments, and cartilage. Imaging evaluates fractures, arthritis, tumors, infection, soft tissue masses, and sports injuries. The choice of modality depends on the structure of interest and the clinical question.

\section{When is MSK imaging ordered?}
\begin{itemize}
    \item \textbf{Suspected fracture:} After trauma---X-ray is first-line.
    \item \textbf{Joint pain:} Assessment of arthritis, internal derangement, or effusion.
    \item \textbf{Suspected bone tumor:} A painful bone lesion or pathological fracture.
    \item \textbf{Tendon or ligament injury:} Suspected rotator cuff tear, ACL rupture, Achilles tendon tear.
    \item \textbf{Suspected infection:} Osteomyelitis or septic arthritis.
    \item \textbf{Soft tissue mass:} Characterising a lump.
\end{itemize}

\section{Modality selection}

\begin{table}[h]
\centering
\caption{MSK imaging modalities}
\label{tab:msk_modality}
\begin{tabular}{@{}p{0.14\textwidth}p{0.39\textwidth}p{0.33\textwidth}@{}}
\hline
\textbf{Modality} & \textbf{Best for} & \textbf{Limitations} \\
\hline
CT & Complex fractures, pre-surgical planning, bone detail & Radiation; poor soft tissue contrast \\
MRI & Marrow, cartilage, tendons, ligaments, infection, tumor extent & Cost, time, contraindications (pacemakers) \\
Ultrasound & Tendons, ligaments, soft tissue masses, dynamic assessment & Operator dependent; cannot see through bone \\
X-ray & Fractures, alignment, arthritis, bone tumors & Low sensitivity for stress fractures, soft tissue \\
\hline
\end{tabular}
\end{table}

\section{Fracture classification}

\subsection{Salter-Harris classification (pediatric growth plate injuries)}

\begin{table}[h]
\centering
\caption{Salter-Harris classification}
\label{tab:salter_harris}
\begin{tabular}{@{}p{0.14\textwidth}p{0.34\textwidth}p{0.12\textwidth}p{0.25\textwidth}@{}}
\hline
\textbf{Type} & \textbf{Description} & \textbf{Frequency} & \textbf{Prognosis} \\
\hline
I & Through physis only & 6\% & Excellent \\
II & Through physis + metaphysis (most common) & 75\% & Excellent \\
III & Through physis + epiphysis into joint & 8\% & Good (may affect growth) \\
IV & Through metaphysis, physis, and epiphysis & 10\% & Risk of growth disturbance \\
V & Crush injury to physis & 1\% & Poor (growth arrest) \\
\hline
\end{tabular}
\end{table}

\subsection{AO/OTA fracture classification}
The AO classification systemises adult long bone fractures:
\begin{itemize}
    \item \textbf{Location:} Diaphyseal (shaft), metaphyseal, or articular.
    \item \textbf{Morphology:} Simple (two fragments) vs comminuted (multiple fragments).
    \item \textbf{Pattern:} Transverse, oblique, spiral, or segmental.
    \item \textbf{Articular:} Partial (single condyle) or complete (bicondylar/t- or y-shaped).
\end{itemize}

\begin{pearlbox}[Look for the second fracture]
When you see a fracture, always look for a second lesion nearby. Occult fractures (hip in elderly, scaphoid, metatarsal stress) are commonly missed on initial X-ray. If clinical suspicion remains high despite a normal X-ray, proceed to MRI (most sensitive) or CT.
\end{pearlbox}

\section{Arthritis on imaging}

\begin{table}[h]
\centering
\caption{Imaging features of common arthritides}
\label{tab:arthritis}
\begin{tabular}{@{}p{0.20\textwidth}p{0.23\textwidth}p{0.22\textwidth}p{0.19\textwidth}@{}}
\hline
\textbf{Feature} & \textbf{OA} & \textbf{RA} & \textbf{Inflammatory} \\
\hline
Erosions & Marginal (at osteophyte base) & Periarticular, well-defined & Romanus lesions (spine) \\
Joint distribution & Asymmetric, DIP/1st CMC & Symmetric, MCP/wrist/PIP & Sacroiliac, spine (SpA) \\
Joint space & Narrowing (loss of cartilage) & Uniform narrowing & ``Squaring,'' fusion \\
Osteophytes & Prominent & Absent & Absent \\
Periarticular osteopenia & Absent & Early feature & Variable \\
Sacroiliitis & No & No & Bilateral symmetric \\
\hline
\end{tabular}
\end{table}

\section{Bone tumor approach: the mnemonic}
Every bone lesion must be systematically characterized. Use the ``age-location-matrix-zone'' approach:

\begin{enumerate}
    \item \textbf{Age:} Most critical factor. Ewing sarcoma and osteosarcoma peak in adolescence; chondrosarcoma and metastases in older adults; giant cell tumor in 20--40 years.
    \item \textbf{Location:} Medullary vs cortical. Epiphyseal (GCT, chondroblastoma), metaphyseal (osteosarcoma, enchondroma, ABC), diaphyseal (Ewing, lymphoma, fibrous dysplasia).
    \item \textbf{Matrix:} What is the lesion made of? Osteoid (cloudy/bony---osteosarcoma), cartilage (punctate/arc---chondroid), fibrous (ground glass---fibrous dysplasia), fat (radiolucent---lipoma, liposarcoma).
    \item \textbf{Zone of transition:} Narrow zone = benign; wide zone = aggressive/malignant. A permeative pattern (moth-eaten or permeative destruction) indicates malignancy.
\end{enumerate}

\begin{pearlbox}[The most dangerous phrase]
``Observation'' of a suspicious bone lesion without a clear follow-up plan is one of the most dangerous phrases in radiology. Every bone lesion requires either a diagnosis, a clear follow-up protocol, or tissue sampling.
\end{pearlbox}

\section{Tendon injuries}

\subsection{Rotator cuff}
MRI is the gold standard for rotator cuff assessment. Key structures: supraspinatus (most commonly torn), infraspinatus, subscapularis, teres minor.
\begin{itemize}
    \item \textbf{Tendinopathy:} Tendinopathy: increased T2 signal within the tendon without a discrete defect.
    \item \textbf{Partial-thickness tear:} High T2 signal not extending through the full tendon thickness (articular-surface, bursal-surface, or intrasubstance).
    \item \textbf{Full-thickness tear:} Complete discontinuity of the tendon with fluid signal throughout. Retraction and fatty infiltration of the muscle belly indicate chronicity.
\end{itemize}

\subsection{Achilles tendon}
Ultrasound is excellent for Achilles tendon assessment:
\begin{itemize}
    \item \textbf{Tendinopathy:} Fusiform thickening ($>$6~mm), hypoechogenicity, loss of fibrillar pattern.
    \item \textbf{Partial tear:} Hypoechoic defect within the tendon.
    \item \textbf{Complete rupture:} Tendon discontinuity with retraction; ``gap'' filled with haematoma. ``Knee bend test'' on clinical examination confirms loss of plantarflexion.
\end{itemize}

\section{Ligament injuries: ACL}
Anterior cruciate ligament (ACL) tears are among the most common knee ligament injuries. MRI findings:
\begin{itemize}
    \item \textbf{Acute:} Disrupted, wavy, or absent ACL fibres with surrounding edema and hemorrhage. ``Bone bruise'' pattern on the posterolateral tibial plateau and posterolateral femoral condyle (pivot-shift mechanism).
    \item \textbf{Chronic:} Ligamentous attenuation, replacement by scar tissue, or empty intercondylar notch.
    \item \textbf{Associated injuries:} Meniscal tears (lateral meniscus $>$ medial), MCL sprain, bone bruising, lipohemarthrosis.
\end{itemize}

\section{Osteomyelitis}
Osteomyelitis is infection of bone. MRI is the most sensitive imaging modality.

\begin{table}[h]
\centering
\caption{MRI features of osteomyelitis}
\label{tab:osteo_mri}
\begin{tabular}{ll}
\hline
\textbf{Sequence} & \textbf{Finding} \\
\hline
DWI & Restricted diffusion in abscess collections \\
Post-contrast T1 & Enhancement of marrow, periosteum, and soft tissue \\
T1-weighted & Low signal in marrow (replaces normal fat signal) \\
T2/STIR & High signal in marrow (edema) \\
\hline
\end{tabular}
\end{table}

\begin{pearlbox}[MRI is superior to X-ray in osteomyelitis]
X-ray findings of osteomyelitis (periosteal reaction, cortical destruction, lucency) may not appear for 10--14 days after symptom onset. MRI detects marrow edema within 24--48 hours and is the investigation of choice when osteomyelitis is suspected.
\end{pearlbox}

\section{When to recommend MRI vs CT}
\begin{itemize}
    \item \textbf{CT:} Complex fractures, pre-operative fracture planning, calcification, bone detail, occult fractures on X-ray (especially spine and pelvis).
    \item \textbf{MRI:} Marrow pathology (tumor, infection, infarction), cartilage and meniscal injuries, ligament and tendon tears, soft tissue tumor characterisation, stress fractures, and spinal cord evaluation.
\end{itemize}

\section{Reporting checklist}
A structured MSK report should include:
\begin{enumerate}
    \item Bone alignment and joint congruity
    \item Fracture description (location, pattern, intra-articular involvement, displacement)
    \item Growth plate status (in children: Salter-Harris type)
    \item Joint space and articular surfaces
    \item Marrow signal (edema, infiltration, infarction)
    \item Tendons and ligaments (integrity, signal, thickening)
    \item Menisci and cartilage
    \item Soft tissues (effusion, haematoma, mass)
    \item Bone lesion characterisation (if present)
    \item Comparison with prior imaging
\end{enumerate}

\clearpage
